\section{Cloud Micro-Physics} \label{sec:basic_usel_microphys}
%------------------------------------------------------
Cloud microphysics is configured by two namelists.
First, the cloud microphysics component is activated by \nmitem{ATMOS_PHY_MP_DO} in \namelist{PARAM_ATMOS}.
Then, the microphysics scheme and its calling interval are specified in \namelist{PARAM_ATMOS_PHY_MP}.

\editboxtwo{
 \verb|&PARAM_ATMOS | & \\
 \verb|ATMOS_PHY_MP_DO = .true.,| & ; Enable cloud microphysics \\
 \verb|/ | & \\
}
\editboxtwo{
 \verb|&PARAM_ATMOS_PHY_MP | & \\
 \verb|MP_TYPE      = "TOMITA08",| & ; Choose from cloud microphysics schemes in Table \ref{tab:mp_scheme_list} \\
 \verb|TIME_DT      = 60.D0,|      & ; Time interval of cloud microphysics process \\
 \verb|TIME_DT_UNIT = "SEC",|      & ; Unit for \verb|TIME_DT| \\
 \verb|do_precipitation    = .true.,|  & ; Enable sedimentation of precipitation particles \\
 \verb|do_negative_fixer   = .true.,|  & ; Enable negative-value fixer for water substances \\
 \verb|ntmax_sedimentation = 1,|       & ; Number of substeps for sedimentation \\
 \verb|/ | & \\
}

The available choices of \nmitem{MP_TYPE} are listed in Table \ref{tab:mp_scheme_list}.
Several options common to cloud microphysics schemes can be specified in \namelist{PARAM_ATMOS_PHY_MP}.
\nmitem{TIME_DT} and \nmitem{TIME_DT_UNIT} specify the time interval of the cloud microphysics process.
\nmitem{do_precipitation} controls whether sedimentation of precipitation particles is calculated.
When precipitation is calculated, 
\nmitem{do_negative_fixer} controls the fixer for negative values of water substances 
and \nmitem{ntmax_sedimentation} specifies the number of substeps with which sedimentation is calculated.

The Kessler and Tomita08 schemes use the cloud microphysics components provided by the SCALE library.
Therefore, the scheme-dependent parameters are basically the same as those used in SCALE-RM.
For the detailed description of these parameters, 
refer to the SCALE-RM User's Guide.

%%%%%%%%%%%
\begin{table}[t]
\caption{Cloud microphysics schemes available in \scaledg.}
\label{tab:mp_scheme_list}
\begin{center}
\begin{tabularx}{150mm}{lXX}
\hline
\rowcolor[gray]{0.9}
\verb|MP_TYPE| & Description & Reference \\
\hline
\verb|KESSLER|  & Three-class one-moment bulk scheme & \cite{kessler_1969} \\
\verb|TOMITA08| & Six-class one-moment bulk scheme   & \cite{tomita_2008} \\
\verb|LSCOND|   & Large-scale condensation scheme    &  \\
\hline
\end{tabularx}
\end{center}
\end{table}
%%%%%%%%%%%


Note that it is necessary to specify the same scheme for \nmitem{MP_TYPE} in the configuration files for both init and run executions.


%------------------------------------------------------------------------
\subsection{Kessler scheme}
\label{subsec:physics_mp_kessler}
%%
This is a three-class one-moment bulk microphysics scheme based on \cite{kessler_1969}.
It treats warm-rain processes using water vapor, cloud water, and rain water.
Since ice-phase water substances are not included, this scheme is mainly suitable for idealized warm-cloud experiments
or for cases where a simple and computationally inexpensive cloud microphysics scheme is sufficient.
This scheme uses the SCALE library, and its parameters are configured as in SCALE-RM. 

%------------------------------------------------------------------------
\subsection{Tomita08 scheme}
\label{subsec:physics_mp_tomita08}
%%
This is a six-class one-moment bulk microphysics scheme based on \cite{tomita_2008}.
It treats water vapor, cloud water, rain water, cloud ice, snow, and graupel.
This scheme is appropriate for experiments where mixed-phase cloud and precipitation processes should be represented.
This scheme uses the SCALE library, and its parameters are configured as in SCALE-RM. 

%------------------------------------------------------------------------
\subsection{Large-scale condensation scheme}
\label{subsec:physics_mp_lscond}
%%
The large-scale condensation scheme is a simplified moist-physics scheme.
It represents condensation and associated latent heating without explicitly calculating detailed bulk microphysical conversion processes among several hydrometeor categories.
This scheme is useful for idealized moist experiments where the feedback between moisture and atmospheric motion is required, 
but detailed cloud microphysics is not the main target.
The precipitation is not explicitly calculated in this scheme, 
and the condensed water is removed from the atmosphere with a timescale of microphysics timestep.


